% Introduction.
Representational similarity analysis (RSA) is one of the most widely used tools for asking
which artificial networks capture the computations of visual cortex: it summarises each
system by the dissimilarities between its responses to a set of stimuli and compares
these summaries directly, without fitting a mapping \citep{kriegeskorte2008}. A
recurring and surprising result is that networks with random, untrained weights can
rival trained networks as models of visual cortex, in some reports exceeding them
\citep{truzzi2020, kazemian2025}.

A result of this kind was reported in THINGS fMRI \citep{hebart2023}:
an untrained network with uncalibrated batch-normalisation statistics exceeded a
backpropagation-trained network at V1 when images were evaluated at 224\,px, but not at
32\,px, the resolution at which the networks had been trained \citep{evalres}; a related
claim appears in \citet{lrrsa}. Both analyses used a single convention:
globally pooled features, correlation distance between units, and voxel responses
normalised over all trials of the experiment.

Correlation distance across
units is not invariant to constant offsets of individual units over the stimulus set:
when units have different mean responses, the dissimilarity between two stimuli depends
in part on how strongly each stimulus pattern is aligned with the shared profile of unit means. Euclidean distance is invariant to such offsets but not to the norm of
each stimulus's response vector. Different distance measures and normalisations
therefore weight a gain-like quantity and the direction of the response pattern
differently. The role of centring is explicit in related similarity measures
\citep{kornblith2019}, and the choice of dissimilarity measure is known to affect
the reliability of RSA \citep{walther2016, diedrichsen2017}, but its effect on
the ranking of candidate models has rarely been examined.

Here we test this result in four ways: (i) a pre-specified analysis replaces RSA by a fitted
ridge readout on the full feature map; (ii) exploratory analyses, each with predictions and
decision rules committed before it was run, ask where the discrepancy between RSA and ridge
arises, including whether the advantage is carried by the component along the mean profile of
the units; (iii) a multiverse of 36 combinations of unit normalisation, distance measure and
voxel normalisation \citep{steegen2016} tests which conclusions hold across defensible analysis
choices, and six ImageNet architectures test whether they generalise beyond the CIFAR-trained
networks; (iv) an exact decomposition of Euclidean distances separates norm and angular
components.

We find that whether the untrained network beats the trained one at V1 depends on the
analysis. The advantage is not reproduced by the fitted readout, disappears under
correlation distance when model units are standardised, and is absent from the angular
component and from the component orthogonal to the mean profile, whereas raw
correlation-distance RSA closely tracks the mean-profile alignment of model and voxels.
The ranking at LOC favours the trained networks wherever it is resolved. A luminance-and-contrast bank scored above every network under the ridge readout at V1 (point estimates). We
report all pre-stated predictions, including those that failed, and state which earlier claims
\citep{evalres, lrrsa} do not hold.
